% !TEX program = xelatex
% !TEX options = -synctex=1 -interaction=nonstopmode -file-line-error
% example-minimal.tex
% Demonstrates rossbeamerminimal features: epigraph, wide image,
% margin notes, aside, blank slide.
\documentclass[palette=minimal]{rossbeamerminimal}

\title{On the Merits of Restraint}
\subtitle{Designing presentations that breathe}
\author{Eleanor Vance}
\institute{Studio Vance, London}
\date{2026}

\begin{document}

% --- Title ---
\begin{frame}[plain]
  \titlepage
\end{frame}

% --- Epigraph ---
\epigraph{Perfection is achieved, not when there is nothing more to add,
  but when there is nothing left to take away.}{Antoine de Saint-Exup\'ery}

% --- Content with margin note ---
\begin{frame}{less is more}
  Good presentations respect the audience's attention. Every element on
  a slide should earn its place.%
  \marginnote{Edward Tufte calls this the ``data-ink ratio'' --- maximise
    the proportion of ink devoted to data.}

  \vskip12pt
  Three principles:
  \begin{itemize}
    \item Eliminate decoration that does not inform
    \item Use whitespace as a structural element
    \item Let typography do the heavy lifting
  \end{itemize}

  \aside{These ideas draw heavily from Butterick's Practical Typography
    and Tufte's Visual Display of Quantitative Information.}
\end{frame}

% --- Wide image ---
\begin{frame}{visual evidence}
  \wideimage{example-image}{A full-width placeholder image demonstrating
    the \texttt{\textbackslash wideimage} command. In practice, this would
    be a photograph, chart, or diagram.}
\end{frame}

% --- More margin notes ---
\begin{frame}{typography matters}
  The choice of typeface communicates before a single word is read.
  Sans-serif faces (Fira Sans, Inter) work well on screen because
  their open counters remain legible at low resolution.%
  \marginnote{Fira Sans was designed by Erik Spiekermann for Mozilla's
    Firefox OS project.}

  \vskip12pt
  Serif faces (Libertinus, Georgia) carry gravitas and are better
  suited to dense, text-heavy slides --- though they should be used
  sparingly in presentations.

  \aside{When in doubt, choose the typeface that disappears.
    The best typography is invisible.}
\end{frame}

% --- Blank slide (pause / breathing room) ---
\blankslide

% --- Closing ---
\begin{frame}{one last thought}
  \vfill
  The most powerful slide in any deck is the one you choose not to show.

  \vskip18pt
  Design is not about making things look good. It is about making
  things work well.

  \vskip18pt
  {\color{rbTextMuted}eleanor@studiovance.co.uk}
  \vfill
\end{frame}

\end{document}
